\section{Composition and Evaluation}
\label{sec:evaluation}

The cases become more informative when composed. Composition also reveals why
no one primitive should absorb every authority.

\subsection{A project cut across four dimensions}

Kungfu's Project Cut architecture distinguishes the following objects
\cite{projectCut2026}:

\begin{table}[ht]
\centering
\small
\begin{tabular}{p{0.22\linewidth}p{0.25\linewidth}p{0.39\linewidth}}
\toprule
Object & Dimension & Authority \\
\midrule
Git tree/commit & material/publication & Exact bytes, ancestry, and atomic
publication coordinate \\
Xinfa Atlas & spatial/epistemic & Declared knowledge, evidence state, routes,
gaps, and semantic closure at one cut \\
Kungfu Episode & temporal/causal & Actual actions, facts, artifacts, receipts,
proof inputs, and sealed transition evidence \\
Project Cut & cross-object binding & Versioned references among source,
Atlas, Episode delta, policies, omissions, and roots \\
Release Passport & release responsibility & What reviewed state and artifacts
were actually promoted and how that claim was verified \\
\bottomrule
\end{tabular}
\caption{Separate authorities composed without a universal project object.}
\label{tab:composition}
\end{table}

A Project Cut is a binding protocol, not a fourth fact engine. Its dependency
direction is one-way:

\begin{equation}
\begin{split}
\text{source}_n + \mathrm{Atlas}_n &\rightarrow \mathrm{Episodes}_{n+1}\\
\text{source}_{n+1} + \mathrm{Episodes}_{n+1}
  &\rightarrow \mathrm{Atlas}_{n+1}\\
\text{source projection} + \mathrm{Atlas}_{n+1}
  + \mathrm{EpisodeDelta}_{n+1} &\rightarrow \mathrm{ProjectCut}_{n+1}\\
\mathrm{ProjectCut}_{n+1} &\rightarrow \text{Git publication}
  \rightarrow \text{Release Passport}.
\end{split}
\end{equation}

Git ancestry does not become causal truth; Episode does not become complete
project knowledge; Atlas does not become runtime history; Passport does not
decide whether the promoted candidate is valuable. The binding preserves
separate roots and explicit responsibility. This is a concrete KFD-1/2 design:
facts remain non-drifting inside their authorities, and cross-authority trust
starts from inspectable bindings rather than name equality.

The current Project Cut decision is accepted but only partially implemented.
Its schema, provider equivalence, settlement, authority cutover, and full
qualification remain staged. It is used here as a composition design and
falsifiable implementation plan, not as completed evidence.

\subsection{The learning loop}

The composition supports a larger loop:

\begin{equation}
\mathrm{Atlas}_n
\rightarrow \text{guides work}
\rightarrow \mathrm{Episodes}_{n+1}
\rightarrow \text{reality pressure}
\rightarrow \text{assessment}
\rightarrow \mathrm{Atlas}_{n+1}.
\end{equation}

In KFD-5, grounded human judgment and agent reasoning interpret the pressure.
In the KFD-6 experiment, an agent mines many Episodes and proposes a candidate
closure, while held-out evidence and separate promotion authority prevent the
loop from certifying itself. Release Passport records what version of the loop
was actually promoted. The result is not a claim of automatic evolution; it is
an auditable path by which a system can learn without replacing external
consequence with generated narrative.

\subsection{Claim ladder}

Evaluation should separate increasingly strong claims:

\begin{description}[leftmargin=*,style=nextline]
  \item[C0: coherent declaration.]
    The candidate has explicit identity, boundary, authority, lifecycle,
    operations, fact sources, observer, alternatives, and falsifiers.
  \item[C1: reproducible artifact.]
    Independent readers can verify the candidate record, roots, and declared
    evidence from public or exported material.
  \item[C2: operational closure.]
    The object supports its minimum operations under interruption, drift, and
    recovery rather than only validating a schema.
  \item[C3: cross-participant reconstruction.]
    A new human or agent can continue work without private oral context and
    without silently creating a second authority.
  \item[C4: generalization.]
    The candidate reduces reconstruction or coordination cost across projects,
    participants, or providers while preserving local facts and falsifiers.
\end{description}

Historical importance, ecosystem lock-in, or universal applicability lies
beyond C4 and cannot be established by the originating team alone.

\subsection{Evaluation protocols}

\paragraph{Deletion and fuse.}
Remove the candidate from a controlled workflow. Measure which coordination
must be reconstructed, which facts become ambiguous, and whether a narrower
alternative closes the gap. Conversely, test whether the candidate genuinely
fuses previously separate mechanisms or merely renames them.

\paragraph{Cold-start recovery.}
Give a fresh participant only the declared artifacts and roots. Measure time,
private-context requests, unsupported assumptions, and correctness of the
reconstructed next action. Delete caches and projections to test stage-zero
recovery from authoritative material.

\paragraph{Drift and fault injection.}
Change a welded surface without its compatibility boundary; omit an Episode;
present a stale Atlas; alter an artifact after release; or mismatch a receipt
and cut. The system should fail visibly or report degraded/unverifiable rather
than synthesize green state.

\paragraph{Held-out and prospective evidence.}
For KFD-6, separate discovery Episodes from evaluation Episodes. Include cases
where the same internal pattern occurs without the proposed boundary. Run the
candidate in shadow mode, then bounded prospective interventions. Measure
compression, prediction, transfer, disconfirmation, and intervention outcome
against retrieval, summarization, and no-new-primitive baselines.

\paragraph{Independent assessment.}
The candidate generator cannot be the only verifier. A separate evaluator,
policy, evidence cut, or accountable participant must be able to reject the
claim. Promotion authority remains separate from discovery authority and
proportionate to consequence.

\subsection{Current evidence status}

At the July 15, 2026 evidence cut, KFD-1 through KFD-5 are active and KFD-6 is
draft. The KFD package publishes machine-readable decision and schema surfaces.
Xinfa Atlas identity and bounded projections are implemented with retained
qualification references. Buildchain Release Passport operates in stable
releases. Episode has a substantial product and runtime architecture but mixed
implementation status across its identity and portability contracts. Project
Cut is accepted and partial.

This supports C0--C2 unevenly across the cases. It does not yet support a
comparative C3 study, a C4 generalization claim, or activation of KFD-6. Those
gaps define the next research work rather than being hidden by the conceptual
coherence of the model.
